%=======================================================================
%  LABORATORIO N.o 5 - POLARIZACION DEL TRANSISTOR BIPOLAR
%  Curso : Circuitos Electronicos  (2026-2)
%  UCSP  - Facultad de Ingenieria - Ing. Electronica y de Telecomunicaciones
%  Docente     : Ebert San Roman Castillo
%  Laboratorio : Ing. Jeanpier Ancori Sanchez
%
%  Version 1 - Practica nueva, pedida por el docente: polarizacion con
%  divisor de tension y tres puntos Q sobre la misma recta de carga, en este
%  orden: cerca del corte, cerca de la saturacion y en el punto medio. A cada
%  circuito se le inyecta una senal AC para ver como el punto Q limita la
%  excursion de la salida.
%    - Formato del Laboratorio 4 (preambulo, cabecera, encabezado/pie,
%      numeracion en romanos, \cuadrorespuesta, cuestionario de 3 preguntas,
%      conclusiones 5 + 5).
%    - Materiales en el estilo de las guias del docente (Lab. 06 y 07 de
%      2026-1): BC548B o 2N2222A, electroliticos. Resistencias a 1/4 W: la
%      que mas disipa es RC, 44 mW con IC = ICsat = 4,46 mA.
%    - El alumno calcula R2: la guia da solo el V_CE objetivo, las formulas
%      y las tablas vacias. No se imprimen los resultados.
%    - Todo sale de la Analog Discovery Studio: rail fijo de +12 V (0,2 A),
%      W1 y osciloscopio (+-25 V, alcanza para ver el colector en DC).
%
%  Circuito fijo: VCC = 12 V, R1 = 10 kOhm, RC = 2,2 kOhm,
%  RE = RE1 + RE2 = 100 + 390 Ohm, CE = 100 uF en paralelo con RE2,
%  Cin = 10 uF. Con R1 = 47 kOhm el punto cerca de la saturacion variaba
%  de 0,8 a 2,2 V entre beta = 100 y 300 (divisor blando); con 10 kOhm se
%  cumple beta*RE >= 10*R2 y la variacion baja de 0,3 V.
%
%  VALORES DE REFERENCIA PARA EL DOCENTE (beta = 200, VBE = 0,7 V; no se
%  imprimen). ICsat = 4,46 mA, VCE(corte) = 12 V.
%                         Circ. 1 corte   Circ. 2 sat.   Circ. 3 medio
%    VCE objetivo          10,5 V          1,5 V          6 V
%    IC objetivo           0,558 mA        3,90 mA        2,23 mA
%    VB de diseno          0,975 V         2,62 V         1,80 V
%    R2 calculada          886 Ohm         2,86 kOhm      1,78 kOhm
%    R2 comercial (E12)    820 Ohm         2,7 kOhm       1,8 kOhm
%    VTH / RTH             0,909 V/758     2,55 V/2,13 k  1,83 V/1,53 k
%    IB                    2,1 uA          18,4 uA        11,3 uA
%    IC                    0,42 mA         3,68 mA        2,26 mA
%    VB / VE / VC          0,91/0,21/11,07 2,51/1,81/3,91 1,81/1,11/7,03
%    VCE                   10,86 V         2,09 V         5,91 V
%    Av = -RC/(re+RE1)     -13,6           -20,6          -19,7
%    vin max (corte)       68 mV pico      394 mV pico    252 mV pico
%    vin max (saturacion)  748 mV pico     88 mV pico     277 mV pico
%  Con 150 mV pico: circuito 1 recorta el semiciclo SUPERIOR de vC (corte),
%  circuito 2 el INFERIOR (saturacion) y circuito 3 sale limpio; empieza a
%  recortar hacia 250 mV pico. R2 = 3,0 kOhm (E24) acerca el circuito 2 a
%  VCE = 0,95 V, pero con tolerancia del 5 % puede quedar saturado.
%
%  SIMULACION (ngspice 36, modelos BC548B de NXP y 2N2222A; netlists y
%  graficas en ../Simulaciones):
%                         Circ. 1 corte   Circ. 2 sat.   Circ. 3 medio
%    VB / VE / VC          0,91/0,29/10,71 2,52/1,85/3,73 1,82/1,16/6,82
%    VCE / IC (BC548B)     10,43 V/0,58 mA 1,88 V/3,76 mA 5,66 V/2,36 mA
%    VCE / IC (2N2222A)    10,48 V/0,57 mA 2,03 V/3,70 mA 5,75 V/2,32 mA
%    beta en el punto Q    328             283            303
%    Av (20 mV pico)       -15,0           -20,2          -19,5
%    Con 150 mV pico       vC tope plano   vC piso plano  limpia
%                          en 12 V (THD    en 1,95 V (THD (THD 1,6 %)
%                          21 %)           9 %)
%    Empieza a recortar    60-80 mV (se    ~90 mV         ~250-300 mV
%                          comprime arriba,
%                          gradual)
%  El circuito 1 se aparta mas del calculo a mano: con 0,5 mA el VBE real
%  es ~0,62 V y no 0,7 V, asi que IC sale ~40 % mayor. Sirve para discutir
%  la aproximacion VBE = 0,7 V.
%=======================================================================

\documentclass[11pt,a4paper]{article}

%------------------------------------------------------------- idioma
% Compilar con XeLaTeX (no con pdflatex): Trebuchet MS es una fuente TrueType
% del sistema y solo fontspec, bajo XeLaTeX o LuaLaTeX, sabe cargarla.
\usepackage{fontspec}
\setmainfont{Trebuchet MS}[
  Path            = C:/Windows/Fonts/,
  Extension       = .ttf,
  UprightFont     = trebuc,
  BoldFont        = trebucbd,
  ItalicFont      = trebucit,
  BoldItalicFont  = trebucbi,
  Ligatures       = TeX
]
\usepackage[spanish,es-nodecimaldot]{babel}
\usepackage{microtype}

%----------------------------------------------------------- geometria
\usepackage[left=2.4cm,right=2.4cm,top=3.1cm,bottom=2.6cm,headheight=30pt,
            footskip=1.2cm]{geometry}

%------------------------------------------------------------ graficos
\usepackage{graphicx}
\usepackage{xcolor}
\usepackage{float}
\usepackage{placeins}
\usepackage[font=small,labelfont=bf,labelsep=period]{caption}
\usepackage{subcaption}
\graphicspath{{imagenes/}}
\usepackage[american,siunitx]{circuitikz}
% babel-spanish activa < > " y rompe las opciones de TikZ sin esta libreria.
\usetikzlibrary{calc,patterns,babel}

%-------------------------------------------------------------- tablas
\usepackage{array}
\usepackage{tabularx}
\usepackage{multirow}

%----------------------------------------------------- matematicas / SI
\usepackage{amsmath,amssymb}
\usepackage{siunitx}
\sisetup{
  output-decimal-marker = {,},
  per-mode              = symbol,
  mode                  = match,
  detect-weight         = true,
  detect-family         = true
}

%--------------------------------------------------------------- listas
\usepackage{enumitem}
\setlist[itemize]{itemsep=3pt,topsep=4pt,parsep=0pt}
\setlist[enumerate]{itemsep=8pt,topsep=4pt,parsep=0pt}

%------------------------------------------------------------ comandos
\usepackage{pgffor}

% Lineas punteadas para respuestas del alumno: \lineasresp{n}
\newcommand{\lineasresp}[1]{%
  \par\vspace{6pt}%
  \foreach \i in {1,...,#1}{\noindent\rule{\linewidth}{0.4pt}\par\vspace{8pt}}%
}

% Recuadro "Figura / Comentario": una celda alta para pegar la captura del
% osciloscopio y otra debajo para el comentario.
\newcommand{\cuadrorespuesta}{%
  \par\vspace{6pt}\noindent
  \begin{tabularx}{\linewidth}{|>{\centering\arraybackslash}X|}
    \hline
    \rule{0pt}{4.4cm}\raisebox{2.1cm}{\textcolor{black!45}{Figura}}\\
    \hline
    \rule{0pt}{2.0cm}\raisebox{0.9cm}{\textcolor{black!45}{Comentario}}\\
    \hline
  \end{tabularx}\par\vspace{8pt}%
}

% Indicacion de la captura que sustenta la fila Ganancia, seguida del
% recuadro Figura / Comentario. \needspace evita que el texto y el recuadro
% queden en paginas distintas.
\newcommand{\capturaganancia}{%
  \par\needspace{8.2cm}\vspace{4pt}\noindent
  \textbf{Captura de la ganancia.} Con una entrada que no produzca recorte, muestre en la
  misma pantalla del osciloscopio la señal de entrada $v_{in}$ (CH1) y la de salida $v_C$
  (CH2). Coloque \textbf{marcadores} (cursores) en el pico máximo y en el pico mínimo de
  cada señal para leer sus amplitudes pico a pico, y calcule
  $A_v = v_{C,pp}/v_{in,pp}$. La captura debe mostrar los marcadores y sus valores; anote el
  resultado en la columna \emph{Medido} de la fila \emph{Ganancia}.%
  \cuadrorespuesta
}

% Celda vacia de altura comoda para escribir a mano.
\newcommand{\hueco}{\rule{0pt}{0.5cm}}

% Tabla de mediciones DC de un circuito: Teorico / Simulado / Medido.
\newcommand{\tablamedicion}[2]{%
  \begin{table}[H]
    \centering
    \renewcommand{\arraystretch}{1.15}
    \begin{tabularx}{0.85\linewidth}{|>{\centering\arraybackslash}m{3.2cm}|*{3}{>{\centering\arraybackslash}X|}}
      \hline
      \textbf{Parámetro} & \textbf{Teórico} & \textbf{Simulado} & \textbf{Medido}\\
      \hline
      $R_2$ (comercial) & \multicolumn{3}{c|}{\hueco} \\ \hline
      $V_B$      & \hueco & & \\ \hline
      $V_E$      & \hueco & & \\ \hline
      $V_C$      & \hueco & & \\ \hline
      $V_{CE}$   & \hueco & & \\ \hline
      $I_B$      & \hueco & & \\ \hline
      $I_C$      & \hueco & & \\ \hline
      $I_E$      & \hueco & & \\ \hline
      $\beta = I_C/I_B$ & \hueco & & \\ \hline
      Ganancia $A_v = v_C/v_{in}$ & \hueco & & \\ \hline
    \end{tabularx}
    \caption{#1}
    \label{#2}
  \end{table}%
}

%----------------------------------------------------- compaginacion
\usepackage{needspace}
\usepackage{etoolbox}
\pretocmd{\section}{\needspace{5\baselineskip}}{}{}
\pretocmd{\subsection}{\needspace{4\baselineskip}}{}{}
\widowpenalty=10000
\clubpenalty=10000

%---------------------------------------------------- encabezado / pie
\usepackage{fancyhdr}
\pagestyle{fancy}
\fancyhf{}
\fancyhead[L]{\footnotesize UCSP --- Facultad de Ingeniería\\
              Ing. Electrónica y de Telecomunicaciones}
\fancyhead[R]{\footnotesize 2026-2 Circuitos Electrónicos\\
              Ebert San Román Castillo}
\fancyfoot[L]{\footnotesize Laboratorio N.\textsuperscript{o} 5 --- Polarización del transistor bipolar}
\fancyfoot[R]{\footnotesize Página \thepage\ de \pageref{LastPage}}
\renewcommand{\headrulewidth}{0.4pt}
\renewcommand{\footrulewidth}{0.4pt}
\usepackage{lastpage}

%------------------------------------------------------- numeracion
\renewcommand{\thesection}{\Roman{section}}
\renewcommand{\thesubsection}{\arabic{subsection}}
\renewcommand{\thesubsubsection}{\thesubsection.\arabic{subsubsection}}
\setcounter{secnumdepth}{3}

%-------------------------------------------------------- hiperenlaces
\usepackage[hidelinks,pdfusetitle]{hyperref}
\usepackage{xurl}
\hypersetup{
  pdftitle={Laboratorio 5 - Polarizacion del transistor bipolar},
  pdfsubject={Circuitos Electronicos - Guia de practicas},
  pdfauthor={Ebert San Roman Castillo; Jeanpier Ancori Sanchez},
  pdfkeywords={electronica, transistor, BJT, polarizacion, recta de carga,
               punto Q, corte, saturacion, divisor de tension,
               Analog Discovery Studio, Proteus, Multisim, osciloscopio}
}

% babel-spanish rotula las tablas como "Cuadro"; los originales dicen "Tabla".
\addto\captionsspanish{\renewcommand{\tablename}{Tabla}}

\begin{document}

%=======================================================================
%  CABECERA DE LA GUIA
%=======================================================================
\begin{center}
    {\normalsize\bfseries SEGUNDA UNIDAD: TRANSISTOR DE UNIÓN BIPOLAR (BJT)}
\end{center}

\vspace{0.55cm}

\begingroup
\renewcommand{\tabularxcolumn}[1]{m{#1}}
\renewcommand{\arraystretch}{2.0}
\setlength{\tabcolsep}{12pt}
\noindent
\begin{tabularx}{\textwidth}{|>{\centering\arraybackslash}m{2.2cm}|>{\centering\arraybackslash}X|>{\centering\arraybackslash}m{3.0cm}|}
\hline
\rule{0pt}{2.1cm}{\fontsize{58}{58}\selectfont\bfseries 5}
    & \begin{minipage}[c]{\linewidth}\centering
          {\large\bfseries Guía de Prácticas}\\[7pt]
          {\Large\bfseries Polarización del\\ Transistor Bipolar}
      \end{minipage}
    & {\Large\bfseries Nota}\\
\hline
\multicolumn{3}{|l|}{\rule{0pt}{1.3cm}\large\textbf{Grupo:}~\rule{4.6cm}{0.4pt}%
    \hfill\large\textbf{Fecha:}~\rule{4.6cm}{0.4pt}\hspace{2mm}}\\
\hline
\multicolumn{3}{|l|}{\rule{0pt}{1.2cm}\large\textbf{Alumno(s):}~\rule{11.2cm}{0.4pt}\hspace{2mm}}\\
\multicolumn{3}{|l|}{\rule{0pt}{1.2cm}\phantom{\large\textbf{Alumno(s):}}~\rule{11.2cm}{0.4pt}\hspace{2mm}}\\
\hline
\end{tabularx}
\endgroup

\vspace{0.8cm}

%=======================================================================
\section{Objetivos}
%=======================================================================
\begin{itemize}
    \item Diseñar la polarización por divisor de tensión de un transistor bipolar para
          ubicar su punto de operación (punto Q) en una posición determinada de la recta de
          carga.
    \item Trazar la recta de carga DC y ubicar sobre ella los puntos Q teóricos, simulados
          y medidos.
    \item Comprobar experimentalmente cómo la posición del punto Q ---cerca del corte, cerca
          de la saturación o en el punto medio--- limita la excursión de una señal AC.
\end{itemize}

%=======================================================================
\section{Contenido teórico}
%=======================================================================
\begin{itemize}
    \item Transistor de unión bipolar (BJT): regiones de trabajo (corte, activa y
          saturación).
    \item Polarización por divisor de tensión, recta de carga y punto de operación.
    \item Amplificador en emisor común: ganancia de voltaje y máxima excursión de salida.
\end{itemize}

%=======================================================================
\section{Equipos y materiales}
%=======================================================================
\noindent\textbf{Laboratorio:} Electrónica de Comunicaciones.

\vspace{0.25cm}
\noindent\textbf{Equipos y dispositivos:}
\begin{itemize}
    \item Tarjeta Analog Discovery Studio: riel fijo de \qty{+12}{\volt}, generador de
          funciones \texttt{W1} y osciloscopio de dos canales.
\end{itemize}

\noindent\textbf{Software:}
\begin{itemize}
    \item Simulador Multisim o Proteus.
\end{itemize}

\vspace{0.25cm}
\needspace{6\baselineskip}
\noindent\textbf{Materiales y fungibles} (deben ser adquiridos por los alumnos):
\begin{itemize}
    % --- Transistores
    \item 02 transistores de unión bipolar NPN BC548B o 2N2222A (uno de repuesto).
    % --- Resistencias
    \item $R_1$: 02 resistencias de \qty{10}{\kilo\ohm} a \qty{0.25}{\watt}.
    \item $R_C$: 02 resistencias de \qty{2.2}{\kilo\ohm} a \qty{0.25}{\watt}.
    \item $R_{E1}$: 02 resistencias de \qty{100}{\ohm} a \qty{0.25}{\watt}.
    \item $R_{E2}$: 02 resistencias de \qty{390}{\ohm} a \qty{0.25}{\watt}.
    % R2 sale 885-890 Ohm, 2,83-2,93 kOhm y 1,77-1,81 kOhm para beta de 100 a
    % 300 (E12: 820 Ohm, 2,7 kOhm, 1,8 kOhm), asi que se piden esos tres.
    \item $R_2$: 02 resistencias de \qty{820}{\ohm}, 02 de \qty{2.7}{\kilo\ohm} y 02 de
          \qty{1.8}{\kilo\ohm}, a \qty{0.25}{\watt}. Se recomienda traer además algunos
          valores vecinos (\qty{680}{\ohm}, \qty{1}{\kilo\ohm}, \qty{1.5}{\kilo\ohm},
          \qty{2.2}{\kilo\ohm}, \qty{3.3}{\kilo\ohm}) por si su diseño previo da un valor
          distinto.
    % --- Capacitores
    \item $C_{in}$: 02 capacitores electrolíticos de \qty{10}{\micro\farad}, de
          \qty{16}{\volt} como mínimo.
    \item $C_E$: 02 capacitores electrolíticos de \qty{100}{\micro\farad}, de
          \qty{16}{\volt} como mínimo.
    % --- Montaje
    \item Protoboard.
    \item Jumpers macho-macho para protoboard.
    \item Alicates para cortar y pelar alambre.
    \item Multímetro (con escala de \unit{\micro\ampere} y, si es posible, medición de
          $h_{FE}$).
    \item Lentes de protección y mandil.
\end{itemize}

\FloatBarrier
%=======================================================================
\section{Conocimientos previos}
%=======================================================================

%-----------------------------------------------------------------------
\subsection{Recta de carga y punto de operación}
%-----------------------------------------------------------------------
Al recorrer la malla colector-emisor del circuito de la figura~\ref{fig:circuito} se obtiene
la ecuación de la \textbf{recta de carga DC}, donde $R_E = R_{E1}+R_{E2}$:
\begin{equation}
  V_{CE} = V_{CC} - I_C\,(R_C + R_E)
  \label{eq:recta}
\end{equation}
Sus dos extremos son el punto de \textbf{saturación} ($V_{CE}\approx 0$) y el de
\textbf{corte} ($I_C = 0$):
\begin{equation}
  I_{C(\text{sat})} = \frac{V_{CC}}{R_C + R_E}
  \qquad\qquad
  V_{CE(\text{corte})} = V_{CC}
  \label{eq:extremos}
\end{equation}
El transistor siempre opera en algún punto de esta recta: el \textbf{punto Q}
$(V_{CEQ},\,I_{CQ})$. La polarización decide dónde queda
(figura~\ref{fig:recta-concepto}).

\begin{figure}[H]
  \centering
  \begin{tikzpicture}[x=0.85cm,y=1.35cm,font=\small]
    % zonas
    \fill[red!12]  (0,0) rectangle (0.25,4.3);
    \fill[blue!12] (0,0) rectangle (11.4,0.08);
    % ejes
    \draw[->] (0,0) -- (11.6,0) node[above] {$V_{CE}$};
    \draw[->] (0,0) -- (0,4.4) node[left] {$I_C$};
    % recta
    \draw[very thick] (0,3.6) -- (10.4,0);
    \node[left]  at (0,3.6) {$\dfrac{V_{CC}}{R_C+R_E}$};
    \node[below] at (10.4,0) {$V_{CC}$};
    \node[below] at (0,0) {0};
    % puntos Q
    \foreach \x/\c/\t/\a in {9.3/blue/{Cerca del corte}/above right,
                            5.2/green!55!black/{Punto medio}/above right,
                            1.4/red/{Cerca de la saturación}/above right}{
      \pgfmathsetmacro{\y}{3.6*(1-\x/10.4)}
      \draw[dotted,\c] (\x,0) -- (\x,\y) -- (0,\y);
      \fill[\c] (\x,\y) circle (3.2pt);
      \node[\a,\c] at (\x,\y) {\t};
    }
    \node[red!70!black,rotate=90,font=\scriptsize] at (0.55,1.3) {saturación};
    \node[blue!70!black,font=\scriptsize] at (8.4,0.25) {corte};
  \end{tikzpicture}
  \caption{Recta de carga DC y tres posiciones posibles del punto Q.}
  \label{fig:recta-concepto}
\end{figure}

%-----------------------------------------------------------------------
\subsection{Polarización por divisor de tensión}
%-----------------------------------------------------------------------
Visto desde la base, el divisor $R_1$--$R_2$ equivale a una fuente de Thévenin:
\begin{equation}
  V_{TH} = V_{CC}\,\frac{R_2}{R_1 + R_2}
  \qquad\qquad
  R_{TH} = \frac{R_1\,R_2}{R_1 + R_2}
  \label{eq:thevenin}
\end{equation}
y la malla base-emisor da la corriente de base:
\begin{equation}
  I_B = \frac{V_{TH} - V_{BE}}{R_{TH} + (\beta+1)\,R_E},
  \qquad
  I_C = \beta\,I_B,
  \qquad
  I_E = I_C + I_B
  \label{eq:ib}
\end{equation}
Con ellas se obtienen los voltajes de los terminales, medidos respecto a tierra:
\begin{equation}
  V_E = I_E\,R_E,
  \qquad
  V_B = V_E + V_{BE},
  \qquad
  V_C = V_{CC} - I_C\,R_C,
  \qquad
  V_{CE} = V_C - V_E
  \label{eq:voltajes}
\end{equation}
El punto Q casi no depende de $\beta$ cuando el divisor es \textbf{firme}:
\begin{equation}
  \beta\,R_E \;\geq\; 10\,R_2
  \label{eq:firme}
\end{equation}

\needspace{10\baselineskip}
\noindent\textbf{Procedimiento de diseño.} Dados $V_{CC}$, $R_1$, $R_C$, $R_E$ y el
$V_{CE}$ que se quiere:
\begin{enumerate}[label=\arabic*), leftmargin=*, itemsep=3pt]
  \item $I_C = \dfrac{V_{CC} - V_{CE}}{R_C + R_E}$ \quad (de la recta de
        carga~\eqref{eq:recta}).
  \item $I_B = \dfrac{I_C}{\beta}$, \quad $I_E = I_C + I_B$.
  \item $V_E = I_E\,R_E$, \quad $V_B = V_E + \qty{0.7}{\volt}$.
  \item Corriente por $R_1$: \quad $I_1 = \dfrac{V_{CC} - V_B}{R_1}$. \quad Por $R_2$ circula
        $I_2 = I_1 - I_B$.
  \item $R_2 = \dfrac{V_B}{I_1 - I_B}$. \quad Se elige el valor comercial más cercano y se
        \textbf{recalcula} el punto Q con \eqref{eq:thevenin}--\eqref{eq:voltajes}.
\end{enumerate}

%-----------------------------------------------------------------------
\subsection{El punto Q y la señal AC}
%-----------------------------------------------------------------------
Con $C_E$ en paralelo con $R_{E2}$, en AC solo queda $R_{E1}$ en el emisor. La ganancia de
voltaje del emisor común es aproximadamente:
\begin{equation}
  A_v = \frac{v_c}{v_{in}} \approx -\frac{R_C}{r_e + R_{E1}},
  \qquad
  r_e = \frac{\qty{26}{\milli\volt}}{I_E}
  \label{eq:ganancia}
\end{equation}
El signo negativo indica que la salida está invertida. La señal se mueve alrededor del
punto Q, y el punto Q fija cuánto puede moverse antes de deformarse:
\begin{itemize}
  \item Hacia el \textbf{corte}, $i_c$ no puede bajar de cero: el colector sube como máximo
        $\Delta v_C \approx I_{CQ}\,R_C$ y se \textbf{recorta el semiciclo superior} de
        $v_C$.
  \item Hacia la \textbf{saturación}, $V_{CE}$ no puede bajar de
        $V_{CE(\text{sat})}\approx\qty{0.2}{\volt}$: el colector baja como máximo
        $\Delta v_C \approx (V_{CEQ} - 0.2)\,\dfrac{R_C}{R_C + R_{E1}}$ y se
        \textbf{recorta el semiciclo inferior}.
\end{itemize}
La amplitud máxima de entrada sin recorte es la menor de esas dos excursiones dividida entre
$|A_v|$:
\begin{equation}
  v_{in,\max} = \frac{\min\!\left(\Delta v_{C,\text{corte}},\;\Delta v_{C,\text{sat}}\right)}{|A_v|}
  \label{eq:vinmax}
\end{equation}

\FloatBarrier
\needspace{14\baselineskip}
%=======================================================================
\section{Desarrollo del laboratorio}
%=======================================================================
Se trabaja siempre con el circuito de la figura~\ref{fig:circuito}. Entre un experimento y otro
\textbf{solo cambia $R_2$}: la recta de carga es la misma y el punto Q se desplaza sobre
ella.

\begin{figure}[H]
  \centering
  \begin{circuitikz}[scale=0.95, transform shape]
    \coordinate (top) at (0,5);
    \coordinate (gnd) at (0,-3.6);
    % transistor
    \draw (4.2,0.6) node[npn](Q){};
    \node[right] at (Q.east) {BC548B};
    % base y divisor
    \coordinate (b) at (2.6,0.6);
    \draw (b) -- (Q.B);
    \draw (b |- top) to[R, l=$R_1$, a=\qty{10}{\kilo\ohm}] (b);
    \draw (b) to[R, l=$R_2$, a={diseño}, *-] (b |- gnd) node[ground]{};
    % colector
    \draw (Q.C) -- ++(0,0.6) coordinate(c) to[R, l=$R_C$, a=\qty{2.2}{\kilo\ohm}] (c |- top);
    \draw (c) to[short, *-o] ++(1.6,0) node[right]{CH2 ($v_C$)};
    % riel
    \draw (b |- top) -- (c |- top) node[circ]{} -- ++(1.4,0) node[right]{$V_{CC} = \qty{+12}{\volt}$};
    % emisor
    \draw (Q.E) to[R, l=$R_{E1}$, a=\qty{100}{\ohm}] ++(0,-1.6) coordinate(e2);
    \draw (e2) to[R, l=$R_{E2}$, a=\qty{390}{\ohm}, *-] (e2 |- gnd) node[ground]{};
    \draw (e2) -- ++(1.9,0) coordinate(ce)
          to[eC, l={$C_E = \qty{100}{\micro\farad}$}] (ce |- gnd) -- (e2 |- gnd);
    % entrada
    \draw (-1.2,-3.6) node[ground]{} to[sV, l=\texttt{W1}] (-1.2,0.6)
          to[short, -*] (-0.6,0.6) node[above]{CH1}
          to[eC, l=$C_{in}$, a=\qty{10}{\micro\farad}, invert] (b);
  \end{circuitikz}
  \caption{Emisor común con polarización por divisor de tensión. Alimentación:
           riel de \qty{+12}{\volt} de la tarjeta. El terminal positivo de los
           electrolíticos va hacia la base ($C_{in}$) y hacia el emisor ($C_E$).}
  \label{fig:circuito}
\end{figure}

% Diagrama de conexion con la tarjeta. Mismo circuito que la figura
% anterior, con W1 como fuente y los cables de la Analog Discovery Studio
% en los colores de su arnes (W1 amarillo, osciloscopio 1 naranja y 2 azul,
% GND negro, alimentacion roja). Los canales son diferenciales: 1- y 2- van
% a GND o la medida flota.
\begin{figure}[H]
  \centering
  \begin{circuitikz}[scale=0.88, transform shape]
    \coordinate (top) at (0,5);
    \coordinate (gnd) at (0,-3.6);
    \colorlet{cw1}{yellow!70!black}
    \colorlet{cch1}{orange!90!black}
    \colorlet{cch2}{blue!75!black}
    \colorlet{cvcc}{red!80!black}
    % transistor y divisor (igual que la figura del circuito)
    \draw (4.2,0.6) node[npn](Q){};
    \node[right] at (Q.east) {BC548B};
    \coordinate (b) at (2.6,0.6);
    \draw (b) -- (Q.B);
    \draw (b |- top) to[R, l=$R_1$, a=\qty{10}{\kilo\ohm}] (b);
    \draw (b) to[R, l=$R_2$, a={diseño}, *-*] (b |- gnd);
    \draw (Q.C) -- ++(0,0.6) coordinate(c) to[R, l=$R_C$, a=\qty{2.2}{\kilo\ohm}] (c |- top);
    \draw (b |- top) -- (c |- top) node[circ]{};
    \draw (Q.E) to[R, l=$R_{E1}$, a=\qty{100}{\ohm}] ++(0,-1.6) coordinate(e2);
    \draw (e2) to[R, l=$R_{E2}$, a=\qty{390}{\ohm}, *-*] (e2 |- gnd);
    \draw (e2) -- ++(1.9,0) coordinate(ce)
          to[eC, l={$C_E = \qty{100}{\micro\farad}$}] (ce |- gnd) node[circ]{};
    \draw (-0.6,0.6) to[eC, l=$C_{in}$, a=\qty{10}{\micro\farad}, invert, *-] (b);
    % bus de tierra
    \coordinate (p2m) at ($(ce)+(1.0,0)$);
    \coordinate (p2p) at ($(ce)+(2.8,0)$);
    \coordinate (p12) at ($(ce)+(3.5,0)$);
    \draw[very thick] (0,-3.6) -- (p2m |- gnd);
    % tarjeta
    \coordinate (ytab) at (0,-5.2);
    \draw[thick, fill=black!6, rounded corners=3pt]
         (-2.0,-5.2) rectangle ($(p12 |- ytab)+(0.9,-1.5)$);
    \node[font=\bfseries] at ($(-2.0,-6.25)!0.5!(p12 |- 0,-6.25)+(0.45,0)$)
         {Analog Discovery Studio};
    \foreach \x/\n/\c in {-1.2/W1/cw1, -0.6/1+/cch1, 0/1$-$/cch1, 0.8/GND/black}{
      \fill[\c] (\x-0.12,-5.32) rectangle (\x+0.12,-5.08);
      \node[font=\scriptsize\ttfamily, below] at (\x,-5.32) {\n};
    }
    \foreach \p/\n/\c in {p2m/2$-$/cch2, p2p/2+/cch2, p12/+12V/cvcc}{
      \fill[\c] ($(\p |- ytab)+(-0.12,-0.12)$) rectangle ($(\p |- ytab)+(0.12,0.12)$);
      \node[font=\scriptsize\ttfamily, below] at ($(\p |- ytab)+(0,-0.12)$) {\n};
    }
    % cables
    \draw[cw1, line width=1.3pt]  (-1.2,-5.08) -- (-1.2,0.6) -- (-0.6,0.6);
    \draw[cch1, line width=1.3pt] (-0.6,-5.08) -- (-0.6,0.6);
    \draw[cch1, line width=1.3pt, dashed] (0,-5.08) -- (0,-3.6) node[circ, black]{};
    \draw[black, line width=1.3pt] (0.8,-5.08) -- (0.8,-3.6) node[circ]{};
    \draw[cch2, line width=1.3pt, dashed] ($(p2m |- ytab)+(0,0.12)$) -- (p2m |- gnd) node[circ, black]{};
    \draw[cch2, line width=1.3pt] ($(p2p |- ytab)+(0,0.12)$) -- (p2p |- c) -- (c) node[circ, black]{};
    \draw[cvcc, line width=1.3pt] ($(p12 |- ytab)+(0,0.12)$) -- (p12 |- top) -- (c |- top);
    % rotulos de las senales
    \node[cch1, font=\scriptsize, above left] at (-0.6,0.6) {CH1: $v_{in}$};
    \node[cch2, font=\scriptsize, above] at ($(c)!0.5!(p2p |- c)$) {CH2: $v_C$};
  \end{circuitikz}
  \caption{Conexión con la Analog Discovery Studio. \texttt{W1} alimenta la entrada;
           \texttt{1+} mide la entrada y \texttt{2+} el colector. Los canales del
           osciloscopio son diferenciales: \texttt{1$-$} y \texttt{2$-$} deben ir a
           \texttt{GND}, o la medida queda flotando. El riel de \qty{+12}{\volt} alimenta el
           circuito y todas las tierras se unen al \texttt{GND} de la tarjeta.}
  \label{fig:conexion}
\end{figure}

%-----------------------------------------------------------------------
\subsection{Diseño previo (traerlo resuelto al laboratorio)}
\label{sec:diseno}
%-----------------------------------------------------------------------
\begin{enumerate}[label=\roman*., leftmargin=*, labelsep=0.6em]
  \item Mida el $\beta$ ($h_{FE}$) de su transistor con el multímetro. Si su multímetro no
        tiene esa función, use $\beta = 200$.
  \item Calcule los extremos de la recta de carga con~\eqref{eq:extremos} y dibújela en la
        figura~\ref{fig:grilla}.
  \item Siguiendo el procedimiento de diseño, calcule $R_2$ para cada uno de los tres
        puntos Q de la tabla~\ref{tab:diseno}. Elija el valor comercial más cercano.
  \item Verifique la condición de divisor firme~\eqref{eq:firme} para cada $R_2$.
  \item Con el valor comercial, recalcule el punto Q, la ganancia~\eqref{eq:ganancia} y la
        amplitud máxima de entrada sin recorte~\eqref{eq:vinmax}. Anote los resultados en la
        columna \emph{Teórico} de las tablas~\ref{tab:circ1}, \ref{tab:circ2}
        y~\ref{tab:circ3} y en la tabla~\ref{tab:ac}.
  \item Simule el circuito de los tres experimentos en Multisim o Proteus y complete la columna
        \emph{Simulado}.
\end{enumerate}

\begin{table}[H]
  \centering
  \renewcommand{\arraystretch}{1.15}
  \begin{tabularx}{0.95\linewidth}{|>{\centering\arraybackslash}m{3.6cm}|*{3}{>{\centering\arraybackslash}X|}}
    \hline
    & \textbf{Experimento 1}\par Cerca del corte
    & \textbf{Experimento 2}\par Cerca de la saturación
    & \textbf{Experimento 3}\par Punto medio \\
    \hline
    $V_{CE}$ objetivo & \qty{10.5}{\volt} & \qty{1.5}{\volt} & $V_{CC}/2 = \qty{6}{\volt}$ \\ \hline
    $I_C$             & \hueco & & \\ \hline
    $I_B$             & \hueco & & \\ \hline
    $I_E$             & \hueco & & \\ \hline
    $V_E$             & \hueco & & \\ \hline
    $V_B$             & \hueco & & \\ \hline
    $I_1$             & \hueco & & \\ \hline
    $R_2$ calculada   & \hueco & & \\ \hline
    $R_2$ comercial   & \hueco & & \\ \hline
    $\beta R_E \geq 10R_2$ ? & \hueco & & \\ \hline
  \end{tabularx}
  \caption{Diseño de $R_2$ para los tres puntos Q ($\beta = $ \rule{1.5cm}{0.4pt}).}
  \label{tab:diseno}
\end{table}

\begin{figure}[H]
  \centering
  \begin{tikzpicture}[x=1.05cm,y=1.45cm,font=\small]
    \draw[step=0.5,black!12,thin] (0,0) grid (13,5);
    \draw[black!30] (0,0) grid (13,5);
    \draw[->] (0,0) -- (13.4,0) node[right] {$V_{CE}$ (V)};
    \draw[->] (0,0) -- (0,5.3)  node[left] {$I_C$ (mA)};
    \foreach \x in {0,1,...,13} \node[below] at (\x,0) {\x};
    \foreach \y in {1,...,5}    \node[left]  at (0,\y) {\y};
  \end{tikzpicture}
  \caption{Recta de carga DC. Dibuje la recta y marque los tres puntos Q: teóricos
           ($\circ$), simulados ($\triangle$) y medidos ($\times$).}
  \label{fig:grilla}
\end{figure}

\FloatBarrier
%-----------------------------------------------------------------------
\subsection{Experimento 1: punto Q cerca del corte}
%-----------------------------------------------------------------------
\begin{enumerate}[label=\roman*., leftmargin=*, labelsep=0.6em]
  \item Arme el circuito de la figura~\ref{fig:circuito} con el $R_2$ de diseño del experimento
        1 y conéctelo a la tarjeta como indica la figura~\ref{fig:conexion}. Alimente con
        el riel de \qty{+12}{\volt} y mida su valor real: si no es exactamente
        \qty{12}{\volt}, recalcule con el valor medido.
  \item \textbf{Polarización (DC).} Con \texttt{W1} apagado, mida $V_B$, $V_E$ y $V_C$ con el
        multímetro, respecto a tierra. Calcule $I_C = (V_{CC} - V_C)/R_C$ e
        $I_E = V_E/R_E$. Mida $I_B$ con el multímetro en serie con la base, en la escala de
        \unit{\micro\ampere}. Complete la tabla~\ref{tab:circ1} y marque el punto Q en la
        figura~\ref{fig:grilla}.
  \item \textbf{Señal AC.} Configure \texttt{W1} como señal sinusoidal de
        \qty{150}{\milli\volt} pico, \qty{1}{\kilo\hertz} y offset \qty{0}{\volt}. Conecte
        CH1 a la salida de \texttt{W1} y CH2 al colector, con \textbf{acoplamiento DC}, para
        ver a la vez el nivel del punto Q y la señal. Capture ambas señales.
  \item ¿Se deforma la salida? ¿Qué semiciclo se recorta y por qué? Reduzca la amplitud de
        \texttt{W1} hasta que la salida quede sin recorte y anote esa amplitud en la
        tabla~\ref{tab:ac}.
\end{enumerate}

\tablamedicion{Experimento 1: punto Q cerca del corte.}{tab:circ1}

\capturaganancia

\clearpage
%-----------------------------------------------------------------------
\subsection{Experimento 2: punto Q cerca de la saturación}
%-----------------------------------------------------------------------
\begin{enumerate}[label=\roman*., leftmargin=*, labelsep=0.6em]
  \item Apague el riel y cambie $R_2$ por la del experimento 2. No modifique nada más.
  \item Repita las mediciones DC del experimento 1 y complete la tabla~\ref{tab:circ2}. Marque
        el punto Q en la figura~\ref{fig:grilla}. Compare el $\beta$ medido con el del
        experimento 1: si es mucho menor, el transistor ya entró en saturación.
  \item Aplique la misma señal de \texttt{W1} (\qty{150}{\milli\volt} pico,
        \qty{1}{\kilo\hertz}) y capture la entrada y la salida.
  \item ¿Qué semiciclo se recorta ahora? ¿Es el mismo que en el experimento 1? Busque la
        amplitud máxima sin recorte y anótela en la tabla~\ref{tab:ac}.
\end{enumerate}

\tablamedicion{Experimento 2: punto Q cerca de la saturación.}{tab:circ2}

\capturaganancia

\clearpage
%-----------------------------------------------------------------------
\subsection{Experimento 3: punto Q en el punto medio}
%-----------------------------------------------------------------------
\begin{enumerate}[label=\roman*., leftmargin=*, labelsep=0.6em]
  \item Cambie $R_2$ por la del experimento 3, repita las mediciones DC y complete la
        tabla~\ref{tab:circ3}. Marque el punto Q en la figura~\ref{fig:grilla}.
  \item Aplique la misma señal de \texttt{W1} (\qty{150}{\milli\volt} pico,
        \qty{1}{\kilo\hertz}) y capture la entrada y la salida. ¿Se recorta?
  \item Aumente poco a poco la amplitud de \texttt{W1} hasta que la salida empiece a
        recortarse. Anote esa amplitud en la tabla~\ref{tab:ac}. ¿Qué semiciclo se recorta
        primero? Capture la señal en ese punto.
\end{enumerate}

\tablamedicion{Experimento 3: punto Q en el punto medio.}{tab:circ3}

\capturaganancia

\FloatBarrier
%-----------------------------------------------------------------------
\subsection{Comparación de los tres puntos Q}
%-----------------------------------------------------------------------
Complete la tabla~\ref{tab:ac}. La ganancia medida es el cociente entre la amplitud pico a
pico de $v_C$ y la de $v_{in}$, con una entrada que no produzca recorte.

\begin{table}[H]
  \centering
  \renewcommand{\arraystretch}{1.15}
  \begin{tabularx}{\linewidth}{|>{\centering\arraybackslash}m{3.1cm}|*{3}{>{\centering\arraybackslash}X|}}
    \hline
    & \textbf{Experimento 1}\par Cerca del corte
    & \textbf{Experimento 2}\par Cerca de la saturación
    & \textbf{Experimento 3}\par Punto medio \\
    \hline
    $V_{CEQ}$ / $I_{CQ}$ medidos          & \hueco & & \\ \hline
    $A_v$ teórica                         & \hueco & & \\ \hline
    $A_v$ medida                          & \hueco & & \\ \hline
    ¿Se recorta con \qty{150}{\milli\volt} pico? & \hueco & & \\ \hline
    Semiciclo de $v_C$ recortado          & \hueco & & \\ \hline
    $v_{in,\max}$ teórica                 & \hueco & & \\ \hline
    $v_{in,\max}$ medida                  & \hueco & & \\ \hline
  \end{tabularx}
  \caption{Punto Q y comportamiento con señal AC.}
  \label{tab:ac}
\end{table}

\FloatBarrier
%=======================================================================
\section{Informe final}
%=======================================================================
\begin{enumerate}[label=\textbf{\arabic*.}, leftmargin=*, labelsep=0.6em]
  \item Adjunte todos los cálculos del diseño previo y la simulación de los tres experimentos.
  \item Presente la recta de carga con los tres puntos Q teóricos, simulados y medidos.
        Explique las diferencias entre ellos.
  \item Presente las formas de onda de cada experimento y explique por qué cada uno se recorta
        de la manera observada.
\end{enumerate}

\FloatBarrier
%=======================================================================
\section{Cuestionario}
%=======================================================================
\begin{enumerate}[label=\textbf{\arabic*.}, leftmargin=*, labelsep=0.6em]
\item Con la misma señal de entrada, el experimento 1 (cerca del corte) y el experimento 2 (cerca de la
saturación) recortan semiciclos distintos de $v_C$. Explique por qué, usando la recta de
carga.
\lineasresp{3}

\item Compare la ganancia medida en los tres experimentos. Si $R_C$ y $R_{E1}$ no cambian,
¿por qué no es la misma? Justifíquelo con la ecuación~\eqref{eq:ganancia}.
\lineasresp{3}

\item Si se reemplaza el BC548B por un 2N2222A, que tiene otro $\beta$, ¿cuánto se movería
el punto Q del experimento 2? ¿Y si $R_1$ fuera de \qty{47}{\kilo\ohm} en lugar de
\qty{10}{\kilo\ohm} (recalculando $R_2$)? Relacione su respuesta con la
condición~\eqref{eq:firme}.
\lineasresp{3}
\end{enumerate}

\FloatBarrier
%=======================================================================
\section{Conclusiones}
%=======================================================================
\begin{itemize}
    \item Redacte como mínimo 05 conclusiones \textbf{cualitativas} y 05
          \textbf{cuantitativas}.
\end{itemize}

\noindent\textbf{Cualitativas}
\lineasresp{10}

\needspace{6\baselineskip}
\noindent\textbf{Cuantitativas}
\lineasresp{10}

\end{document}
